% !TeX root = ../main.tex
\section{Giới Thiệu}
\label{sec:intro}

Các trò chơi cờ trừu tượng đã tồn tại hàng nghìn năm, từ Cờ Vây của
Trung Quốc đến Cờ Đam châu Âu và Reversi ra đời cuối thế kỷ~19.
Mỗi game định nghĩa một không gian chiến thuật riêng biệt qua tập luật
tối giản của nó.

\textbf{Tactical Formation} ra đời từ câu hỏi thiết kế: \emph{điều gì xảy
ra nếu ta kết hợp cơ chế di chuyển của Cờ Đam với cơ chế bắt quân của
Reversi, nhưng không loại bỏ quân bị nhảy qua, và mở rộng sang nhiều
người chơi?}

Sự kết hợp này tạo ra một căng thẳng chiến thuật mới. Trong Cờ Đam,
người chơi nhảy qua quân để \emph{ăn} nó---quân biến mất. Trong Reversi,
người chơi đặt quân mới để \emph{kẹp} quân đối thủ---quân bị lật màu.
Trong Tactical Formation, người chơi nhảy qua một quân bất kỳ (kể cả quân
của chính mình) để di chuyển, rồi kẹp sandwich từ vị trí mới---quân bị
nhảy qua vẫn còn trên bàn và có thể trở thành vũ khí hoặc rào cản ở lượt
tiếp theo.

Kết quả là mỗi nước đi mang hai chiều tác dụng độc lập: vị trí mới của
quân (ảnh hưởng đến di chuyển tương lai) và sandwich tạo ra từ vị trí đó
(ảnh hưởng đến lật quân ngay lập tức). Người chơi phải cân nhắc cả hai
chiều đồng thời, tạo ra chiều sâu chiến thuật khác biệt so với cả hai
nguồn cảm hứng.

Bài báo có các đóng góp sau:
\begin{enumerate}[label=(\roman*)]
  \item Đặc tả đầy đủ, nhất quán luật chơi Tactical Formation cho
        2--4 người chơi (cấu hình xuất phát 3 người hiện chưa cân đối---xem
        Mục~\ref{sec:limitations}).
  \item Mô tả cấu trúc bàn cờ hai lớp (đỉnh và tâm) và phân tích ảnh
        hưởng của nó đến không gian chiến thuật.
  \item Mô hình toán học hình thức cho chế độ 2 người chơi, kèm
        \emph{Định luật Bảo toàn 28 quân} cô đọng tinh thần \emph{ăn = đổi phe}.
  \item Xây dựng AI đồng hành dựa trên Paranoid Minimax, kèm phân tích
        tại sao MCTS không phù hợp cho lớp game này.
  \item Đánh giá thực nghiệm trò chơi và AI qua branching factor (2--4
        người), giải đấu bot-vs-bot (chế độ 2 người) và hiệu năng thời
        gian thực.
  \item Phát hành phần mềm mã nguồn mở để cộng đồng chơi, nghiên cứu
        và phát triển tiếp.
\end{enumerate}

Phần còn lại của bài báo được tổ chức như sau. Mục~\ref{sec:related}
khảo sát các game liên quan. Mục~\ref{sec:rules} đặc tả đầy đủ luật chơi
bằng hình ảnh bàn cờ. Mục~\ref{sec:variants} mô tả các biến thể nhiều
người. Mục~\ref{sec:ai} trình bày thuật toán AI đồng hành.
Mục~\ref{sec:impl} mô tả phần mềm. Mục~\ref{sec:eval} đánh giá thực
nghiệm. Mục~\ref{sec:limitations} nêu hạn chế, và
Mục~\ref{sec:conclusion} kết luận. Mô hình toán học hình thức cho chế độ
2 người được trình bày đầy đủ ở Phụ lục~\ref{app:formal}.
